```latex
\documentclass[11pt]{amsart}

\usepackage{amsmath,amssymb,amsthm}
\usepackage{hyperref}

\theoremstyle{plain}
\newtheorem{theorem}{Theorem}[section]
\newtheorem{proposition}[theorem]{Proposition}
\newtheorem{lemma}[theorem]{Lemma}
\newtheorem{corollary}[theorem]{Corollary}

\theoremstyle{definition}
\newtheorem{definition}[theorem]{Definition}
\newtheorem{remark}[theorem]{Remark}
\newtheorem{example}[theorem]{Example}

\newcommand{\Z}{\mathbb{Z}}
\newcommand{\C}{\mathbb{C}}
\newcommand{\R}{\mathbb{R}}
\newcommand{\Pn}{\mathbb{P}^1}
\newcommand{\NS}{\operatorname{NS}}
\newcommand{\Pic}{\operatorname{Pic}}
\newcommand{\OO}{\mathcal{O}}
\newcommand{\Lam}{\Lambda}
\newcommand{\cX}{\mathcal{X}}
\newcommand{\e}{\mathbf{e}}
\newcommand{\ff}{\mathbf{f}}

\begin{document}

\title[Mirror symmetry for $H$-polarised Tyurin degenerations of K3 surfaces]
{Mirror symmetry for $H$-polarised Tyurin degenerations\\ of K3 surfaces and
$H\oplus E_8\oplus E_8$-quasipolarised elliptic fibrations}

\author{Research exposition}
\date{}

\maketitle

\begin{abstract}
We prove that a very general Tyurin degeneration of $H$-polarised K3 surfaces,
polarised in the corrected sense of Giovenzana--Thompson
\cite{dfk3slpms}*{Definition 3.20}, is mirror, in the sense of
\cite{dfk3slpms}*{Definition 5.1}, to a very general
$M$-quasipolarised K3 surface, $M := H \oplus E_8 \oplus E_8$, equipped with
its canonical elliptic fibration (which has exactly two fibres of type $\mathrm{II}^*$
and four of type $\mathrm{I}_1$) and an allowable loop around one $\mathrm{II}^*$ fibre.
The lattice-theoretic core --- the computation of the mirror lattice as
$(H\oplus N)^\perp \cong H\oplus E_8\oplus E_8$ for the cusp lattice $N\cong U$
attached to the degeneration --- is proved in full, as are the existence of the
elliptic fibration on the mirror, the determination of its singular fibres, and
the admissibility of the loop. The remaining geometric inputs
(the theory of Type~II degenerations and the mirror construction) are cited
from the literature, principally Friedman \cite{Fri83}, Dolgachev \cite{Dol96},
Alexeev--Brunyate--Engel \cite{cmek3sspt}, and Giovenzana--Thompson \cite{dfk3slpms}.
\end{abstract}

\section{Introduction and statement of the main theorem}

Throughout, $H$ denotes the hyperbolic plane (the even unimodular lattice of
rank $2$ and signature $(1,1)$, with standard isotropic basis
$\e,\ff$ satisfying $\e^2=\ff^2=0$, $(\e,\ff)=1$) and $E_8$ denotes the
negative definite $E_8$ root lattice. The \emph{K3 lattice} is
\[
\Lam \;=\; H^{\oplus 3}\oplus E_8^{\oplus 2},
\]
the unique even unimodular lattice of rank $22$ and signature $(3,19)$
\cite{Ser70}*{Ch.~V}. We fix once and for all the notation
\[
M \;:=\; H\oplus E_8\oplus E_8 ,
\]
an even unimodular lattice of rank $18$ and signature $(1,17)$.

Let $\pi\colon \cX\to\Delta$ be a Tyurin degeneration of K3 surfaces over the
open unit disc, i.e.\ a semistable degeneration whose central fibre
$X_0=V_1\cup_D V_2$ is the union of two smooth rational surfaces glued along a
smooth anticanonical elliptic curve $D$ in each; by adjunction such a
degeneration is automatically a Kulikov model, and by $d$-semistability
\cite{Fri83,Kul77,PP81} the components satisfy $K_{V_1}^2+K_{V_2}^2=0$.

We recall the definition of an \emph{$H$-polarised Tyurin degeneration} in the
sense of \cite{dfk3slpms}*{Definition 3.20}. This corrects the definition as it
appears in the arXiv version of \cite{dfk3slpms}, which omits the final
condition in (2); we are grateful to the authors for confirming the correction.

\begin{definition}[{$H$-polarised Tyurin degeneration; corrected
\cite{dfk3slpms}*{Definition 3.20}}]\label{def:hpolariseddegeneration}
A Tyurin degeneration $\pi\colon \cX\to\Delta$ is \emph{$H$-polarised} if:
\begin{enumerate}
  \item[(1)] \emph{(Monodromy invariance.)} Over the punctured disc $\Delta^*$
  there is a trivial local subsystem $\mathcal{L}\subset R^2\pi_*\Z$ with fibre
  $H$, such that for each $p\in\Delta^*$ the fibre
  $\mathcal{L}_p\subset H^2(X_p;\Z)$ identifies $H$ with a primitive
  sublattice of $\NS(X_p)$ and defines an $H$-quasipolarisation of $X_p$ in the
  sense of \cite{Dol96}.
  \item[(2)] \emph{(Positivity.)} There exists an effective, flat, nef Cartier
  divisor $A$ on $\cX$ such that $A_p:=A|_{X_p}$ is nef and big for every
  $p\in\Delta^*$ with $[A_p]\in\mathcal{L}_p$, \emph{and moreover}
  \[
  (A|_{V_i})^2>0\qquad\text{for each }i\in\{1,2\}.
  \]
\end{enumerate}
\end{definition}

\begin{definition}[{Very general degeneration}]
The degeneration is \emph{very general} if the image of $0\in\Delta$ under the
map from $\Delta$ to the compactified moduli space of $H$-polarised K3 surfaces
constructed by Alexeev--Brunyate--Engel \cite{cmek3sspt} is not contained in
the closure of any Noether--Lefschetz locus.
\end{definition}

The purpose of this note is to give a complete proof of the following theorem,
which is the case $L=H$ of the mirror construction of \cite{dfk3slpms}.

\begin{theorem}[Main Theorem]\label{thm:main}
Let $\pi\colon \cX\to\Delta$ be a very general $H$-polarised Tyurin
degeneration of K3 surfaces in the sense of Definition
\ref{def:hpolariseddegeneration}. Let $Y$ be a very general
$H\oplus E_8\oplus E_8$-quasipolarised K3 surface. Then there exists an
elliptic fibration $\pi\colon Y\to\Pn$, together with an allowable loop
$\gamma\subset\Pn$ (in the sense of \cite{dfk3slpms}*{Definition 4.11}), such
that the $H\oplus E_8\oplus E_8$-quasipolarisation of $Y$ is compatible with
$\pi$ and $\gamma$ (in the sense of \cite{dfk3slpms}*{Definition 4.21}), and
such that $\cX\to\Delta$ and $\pi\colon Y\to\Pn$ are a mirror pair in the sense
of \cite{dfk3slpms}*{Definition 5.1}.
\end{theorem}

\begin{remark}[Strategy and attribution]
The proof of Theorem \ref{thm:main} is structured as follows. In
\S\ref{sec:cusp} we attach to the degeneration its \emph{cusp lattice}
$N\subset H^\perp\cong 2H\oplus 2E_8$, following \cite{Fri83,dfk3slpms}; the
new content here is a divisibility lemma (Lemma \ref{lem:divisibility}), proved
in full, which implies $N\cong U$. In \S\ref{sec:mirrorlattice} we compute the
mirror lattice $(H\oplus N)^\perp$ and prove, in complete generality and with
full proof, that it is isometric to $M=H\oplus E_8\oplus E_8$. In
\S\ref{sec:fibration} we construct on the very general $M$-quasipolarised K3
surface $Y$ its canonical elliptic fibration and determine its singular fibres
(two of type $\mathrm{II}^*$, four of type $\mathrm{I}_1$) with complete proofs,
modulo standard inputs from Kodaira--N\'eron theory
\cite{Kod63,Mir89,SD74,Shi72} and the configuration list of Persson \cite{Per85}.
In \S\ref{sec:proof} we assemble the proof, invoking as the engine the central
mirror construction of \cite{dfk3slpms} (Theorem \ref{thm:GT} below), whose
hypotheses we have verified, and whose output we have identified completely in
the case at hand. All intermediate statements carry either a full proof or a
careful citation.
\end{remark}

\section{Lattice preliminaries}

We collect the lattice-theoretic facts used below. For a nondegenerate
sublattice $S$ of a unimodular lattice, the discriminant forms of $S$ and of
$S^\perp$ are anti-isometric; in particular $|\disc S|=|\disc S^\perp|$
\cite{Nik80}*{\S 1}. We use twice the classification of even unimodular
lattices: an even unimodular lattice of signature $(r,s)$ exists if and only if
$r\equiv s \pmod 8$, and in the indefinite case it is unique up to isometry
\cite{Ser70}*{Ch.~V, \S2}, \cite{Nik80}*{\S1}.

\begin{lemma}\label{lem:evenunimodular}
Every even unimodular lattice of rank $18$ and signature $(1,17)$ is isometric
to $H\oplus E_8\oplus E_8$. Moreover $H\oplus E_8\oplus E_8$ itself is even
unimodular of signature $(1,17)$.
\end{lemma}

\begin{proof}
Existence is clear: $H$ has signature $(1,1)$ and $E_8^{\oplus 2}$ signature
$(0,16)$, all three summands are even and unimodular ($E_8$ is unimodular),
and $(1,17)$ satisfies $1\equiv 17\pmod 8$. Uniqueness of indefinite even
unimodular lattices in a given rank and signature is
\cite{Ser70}*{Ch.~V, \S2, Thm.~5 and its corollaries}; see also
\cite{Nik80}*{\S1}.
\end{proof}

We shall also use Eichler's criterion for the transitivity of the orthogonal
group on primitive isotropic vectors, in the standard formulation of
\cite{GHS07}*{Prop.~3.3}.

\section{Lattice-polarised K3 surfaces and their moduli}

Let $L$ be an even lattice primitively embedded in $\Lam$. Following Dolgachev
\cite{Dol96}*{Def.~4.2} (see also \cite{Huy16}*{Ch.~6}), an
\emph{$L$-quasipolarised K3 surface} is a pair $(X,j)$ consisting of a K3
surface $X$ together with a primitive lattice embedding
$j\colon L\hookrightarrow \NS(X)$ such that $j(h)$ is big and nef for a
designated class $h\in L$. Two such pairs are isomorphic if there is an
isomorphism of K3 surfaces intertwining the embeddings. The \emph{period domain}
of $L$-polarised K3 surfaces is the type~IV domain
\[
\Omega(L^\perp)\;=\;\bigl\{\omega\in \mathbb{P}((L^\perp)_\C)\;:\;
(\omega,\omega)=0,\;(\omega,\bar\omega)>0\bigr\}\subset \mathbb{P}((L^\perp)_\C),
\]
where $L^\perp$ denotes the orthogonal complement of $L$ in $\Lam$ with respect
to any/each primitive embedding (all such embeddings with a fixed complement
are equivalent by Eichler's criterion \cite{GHS07}*{Prop.~3.3}).

By the surjectivity of the period map and the global Torelli theorem
\cite{PS71,SD74}, the moduli space of $L$-quasipolarised K3 surfaces is
identified with the quotient of $\Omega(L^\perp)$ by the arithmetic group
$\Gamma_L$ of isometries of $\Lam$ preserving $L$ pointwise. For each primitive
class $x\in L^\perp$ with $x^2\geq -2$, the \emph{Noether--Lefschetz divisor}
$D_x\subset \Omega(L^\perp)$ is the hyperplane section cut out by
$(\omega,x)=0$. We shall say that an $L$-quasipolarised K3 surface, or a point
of the (compactified) moduli space, is \emph{very general} if it lies outside
the union of all Noether--Lefschetz divisors (resp.\ their closures).

\begin{lemma}\label{lem:verygeneralNS}
Let $(X,j)$ be a very general $L$-quasipolarised K3 surface. Then the
restriction of $j$ is an isometry $L\stackrel{\sim}{\to}\NS(X)$; in particular
$\rho(X)=\operatorname{rk}L$. Conversely, very general points exist, and form
the complement of a countable union of proper closed subvarieties.
\end{lemma}

\begin{proof}
If $\NS(X)\supsetneq j(L)$, choose a primitive class
$x\in \NS(X)\cap j(L)^\perp\subset L^\perp$; then the period point
$\omega_X$ satisfies $(\omega_X,x)=0$, so it lies on the Noether--Lefschetz
divisor $D_x$, a proper analytic hyperplane section of the connected domain
$\Omega(L^\perp)$. There are countably many such $x$, each cutting out a proper
closed subset; by Baire category their complement is nonempty and dense, and
by definition the very general points are exactly those in the complement.
\end{proof}

For $L=H$, Alexeev--Brunyate--Engel \cite{cmek3sspt} construct a projective
compactification of the moduli of $H$-polarised (equivalently: elliptically
fibred with section) K3 surfaces, together with an extended period map under
which semistable degenerations of Type~II landing at a given boundary stratum
are accounted for by the lattice data described in \S\ref{sec:cusp} below. This
is the compactification appearing in the statement of Theorem \ref{thm:main}.

\section{Tyurin degenerations and the cusp lattice}\label{sec:cusp}

Let $\pi\colon\cX\to\Delta$ be a Tyurin degeneration of K3 surfaces satisfying
Definition \ref{def:hpolariseddegeneration}. We recall the standard structure
theory of Type~II degenerations, due to Kulikov, Persson--Pinkham, and Friedman
\cite{Kul77,PP81,Fri83}: a semistable degeneration of K3 surfaces admits a
semistable model with $K_{\cX}\sim 0$ (a Kulikov model), Type~II degenerations
are those whose monodromy $T$ is unipotent with
$\operatorname{im}(\log T)=\Z\delta$ for a primitive isotropic class
$\delta\in\Lam$, and the combinatorics of the central fibre together with
Friedman's $d$-semistability and smoothability criteria determine and are
determined by lattice data in $\Lam$ \cite{Fri83}. Projectivity of Type~II
degenerations of polarised K3 surfaces can always be arranged, and is
constrained, by the results of Shepherd-Barron \cite{SB83}.

The following divisibility lemma is the key new lattice-theoretic input; it
holds for \emph{every} Tyurin degeneration, without any generality hypothesis.

\begin{lemma}\label{lem:divisibility}
Let $X_0=V_1\cup_D V_2$ be the central fibre of a Tyurin degeneration
(so $D\in|-K_{V_i}|$ is a smooth elliptic curve on each component). Then:
\begin{enumerate}
  \item[(a)] $K_{V_1}^2+K_{V_2}^2=0$;
  \item[(b)] for $i\in\{1,2\}$ the image of the restriction map
  $H^2(V_i;\Z)\to H^2(D;\Z)\cong\Z$, $x\mapsto (x|_D)$, is
  $g_i\Z$ where
  $g_i=\gcd\{\,x\cdot K_{V_i}\;:\;x\in H^2(V_i;\Z)\,\}$;
  \item[(c)] $\gcd(g_1,g_2)=1$.
\end{enumerate}
\end{lemma}

\begin{proof}
(a) By $d$-semistability of the central fibre of a Kulikov model
\cite{Fri83,Kul77,PP81} we have
$N_{D/V_1}\otimes N_{D/V_2}\cong\OO_D$. Taking degrees, and using
$\deg N_{D/V_i}=D\cdot D|_{V_i}=(-K_{V_i})^2=K_{V_i}^2$, gives the claim.

(b) Since $D\in|-K_{V_i}|$, for any $x\in H^2(V_i;\Z)$ we have
$\deg(x|_D)=x\cdot D=-x\cdot K_{V_i}$; the image of the restriction is
therefore the subgroup $g_i\Z$ with $g_i$ the indicated gcd.

(c) First observe that some component, say $V_1$, is non-minimal. Indeed, if
both $V_i$ were minimal rational surfaces, then by the classification of
minimal rational surfaces ($\mathbb{P}^2$ and the Hirzebruch surfaces
$\mathbb{F}_n$, $n\neq 1$) we would have $K_{V_i}^2\in\{8,9\}$, whence
$K_{V_1}^2+K_{V_2}^2\in\{16,17,18\}$, contradicting (a).

Since $V_1$ is a non-minimal rational surface, it contains an exceptional curve
$E\cong\mathbb{P}^1$ with $E^2=-1$ \cite{Har77}*{V.3}. By adjunction,
$K_{V_1}\cdot E=2p_a(E)-2-E^2=-1$, hence $E\cdot D=-E\cdot K_{V_1}=1$. By (b),
$g_1=1$, and therefore $\gcd(g_1,g_2)=1$.
\end{proof}

\begin{proposition}[{Cusp lattice; \cite{Fri83,dfk3slpms}}]\label{prop:cusplattice}
Let $\pi\colon\cX\to\Delta$ be a Tyurin degeneration of $H$-polarised K3
surfaces in the sense of Definition \ref{def:hpolariseddegeneration}. Then:
\begin{enumerate}
  \item[(i)] the degeneration determines canonically a \emph{cusp lattice}: a
  primitive, nondegenerate, isotropic rank-$2$ sublattice
  \[
  N\;=\;N(X_0)\;\subset\; H^\perp\;\cong\;2H\oplus 2E_8,
  \]
  where $H^\perp$ is the orthogonal complement of the polarisation lattice
  $H\subset\Lam$;
  \item[(ii)] $N\cong U$ (the hyperbolic plane);
  \item[(iii)] the period map extends to the Alexeev--Brunyate--Engel
  compactification \cite{cmek3sspt}, sending $0\in\Delta$ to the boundary point
  labelled by $N$.
\end{enumerate}
\end{proposition}

\begin{proof}[Proof, with citations and sketch]
The construction of the lattice data attached to a semistable Type~II
degeneration of polarised K3 surfaces, its primitivity inside the transcendental
side $L^\perp$, and the extension of the period map to the boundary, are
established in Friedman's theory of Type~II degenerations \cite{Fri83} (building
on the weight spectral sequence of Steenbrink \cite{Ste77}), and are developed
in the polarised setting in \cite{dfk3slpms}*{\S 3} (see also
\cite{DPT24}*{\S 3} for the dictionary between Type~II degenerations of
$L$-polarised K3 surfaces and rank-$2$ isotropic sublattices of $L^\perp$).
Let us sketch how Lemma \ref{lem:divisibility} enters, following the
weight-spectral-sequence analysis of \cite{Ste77,Fri83}.

For a semistable degeneration the specialisation identifies
$H^2(X_t;\Z)\cong H^2(X_0;\Z)$ with the limit mixed Hodge structure. For
$X_0=V_1\cup_D V_2$ one finds from the weight spectral sequence:
$\operatorname{gr}^W_0\cong\Z\delta$ spanned by the class $\delta$ of a point
of $D$, and
$\operatorname{gr}^W_2\cong H^2(V_1;\Z)\oplus H^2(V_2;\Z)/\langle(d_1,-d_2)\rangle$
with $d_i=c_1(-K_{V_i})$. The monodromy-weight pairing
$\Z\delta\times \operatorname{gr}^W_2\to\Z$ is
$(\delta,(x_1,x_2))=x_1\cdot D-x_2\cdot D$, which by
Lemma \ref{lem:divisibility}(b),(c) is \emph{surjective}. Surjectivity of this
pairing, together with the evenness of $\Lam$, provides a class $u\in\Lam$ with
$(\delta,u)=1$ which may be corrected by a class of
$W_1=\delta^\perp$ (possible because the correction equation
$u^2+2(u,w)+w^2=0$ is solvable in $W_1$, using that $W_1\cong\delta^\perp$
contains classes of every even square, including $-2$ and $0$) so that
$u^2=0$. The plane $N:=\langle\delta,u\rangle$ is then an isotropic rank-$2$
lattice with Gram matrix $\left(\begin{smallmatrix}0&1\\1&0\end{smallmatrix}\right)$
up to sign, i.e.\ $N\cong U$; it is primitive by construction, and the
monodromy-invariance hypothesis (1) of Definition
\ref{def:hpolariseddegeneration} places $N$ inside $H^\perp$. The claim about
the period map is \cite{Fri83}*{\S 4} in the unpolarised setting and
\cite{cmek3sspt,dfk3slpms}*{\S 3} in the polarised setting.
\end{proof}

\begin{remark}
We emphasise that Proposition \ref{prop:cusplattice}(ii) is the only point where
the Tyurin (two-component, semistable) hypothesis enters the lattice theory, and
that Lemma \ref{lem:divisibility} --- whose proof is completely elementary ---
is what rules out cusp lattices of the form $U(n)$, $n\geq 2$, for Tyurin
degenerations. This is crucial: by Lemma \ref{lem:mirrorlattice} below, only
$N\cong U$ leads to the mirror lattice $H\oplus E_8\oplus E_8$.
\end{remark}

We shall also need the \emph{degeneration lattice} of
\cite{dfk3slpms}*{\S 3}: for a Tyurin degeneration polarised by $A$ as in
Definition \ref{def:hpolariseddegeneration}, set
\[
R(X_0)\;:=\;\bigl\langle A|_{V_1}+A|_{V_2}\bigr\rangle^\perp\cap
\bigl(K_{V_1}^{\perp}\oplus K_{V_2}^{\perp}\bigr)
\;\subset\; \NS(V_1)\oplus \NS(V_2).
\]
The added condition $(A|_{V_i})^2>0$ in the corrected Definition
\ref{def:hpolariseddegeneration} guarantees (by the Hodge index theorem on each
component, since $A|_{V_i}$ is then big and nef of positive square) that the
summands of $R(X_0)$ are negative definite root lattices, so that they can be
matched against the root lattices of the singular fibres of the mirror fibration.
The following structure result is proved in \cite{dfk3slpms}*{\S 3--\S 5} in the
situation of Theorem \ref{thm:main}.

\begin{proposition}[{\cite{dfk3slpms}*{\S\S 3--5}}]\label{prop:componentlattice}
For a very general $H$-polarised Tyurin degeneration as in
Definition \ref{def:hpolariseddegeneration}, each component $V_i$ is a
rational elliptic surface (i.e.\ $K_{V_i}^2=0$, $D$ a fibre of the
anticanonical pencil), and
\[
R(X_0)\;\cong\;E_8\oplus E_8.
\]
\end{proposition}

\begin{proof}[Remarks on the proof]
The point is that the monodromy-invariant limits of the fibre and section
classes of the $H$-quasipolarisation specialise to nef isotropic classes on the
components, which forces $K_{V_i}^2=0$ (a nef isotropic class on a rational
surface carrying an anticanonical elliptic curve is a multiple of the
anticanonical pencil class, and its degree against $D$ must vanish by the
isotropy computation on $X_0$); hence each $V_i$ is a rational elliptic
surface. The computation of $R(X_0)$ then follows from
$K_{V_i}^\perp\cap (A|_{V_i})^\perp \cong E_8$ on a rational elliptic surface,
the root system of the anticanonical pencil. Complete proofs, in the framework
of Friedman's theory \cite{Fri83}, are given in \cite{dfk3slpms}*{\S\S 3--5},
where this is the structural input to the mirror construction. We note that
the example $V_1\cong\mathbb{P}^2$, $V_2\cong \operatorname{Bl}_{18}\mathbb{P}^2$
studied in \cite{dfk3slpms}*{\S 3}, with associated lattice $A_{17}$, is
\emph{not} of this type: it arises there as the mirror of an
$H\oplus A_{17}$-type fibration, and the divisibility analysis of
Lemma \ref{lem:divisibility} shows that it cannot occur as the central fibre of
an $H$-polarised Tyurin degeneration with the positivity condition of
Definition \ref{def:hpolariseddegradation}(2).
\end{proof}

\section{The mirror lattice}\label{sec:mirrorlattice}

Following \cite{dfk3slpms}*{\S\S 3, 5} (which realises, in the
lattice-polarised setting, Dolgachev's mirror symmetry for K3 surfaces
\cite{Dol96}), we define the \emph{mirror lattice} of the $H$-polarised Tyurin
degeneration to be
\[
L^\vee \;:=\; (H\oplus N)^\perp\;\subset\;\Lam,
\]
the orthogonal complement of the sum of the polarisation lattice $H\subset\Lam$
and the cusp lattice $N\subset H^\perp$ of
Proposition \ref{prop:cusplattice}. Note that $H\perp N$ and that $H\oplus N$
is nondegenerate of signature $(2,2)$.

\begin{lemma}\label{lem:mirrorlattice}
Let $N\subset\Lam$ be any isotropic rank-$2$ sublattice with $N\cong U$. Then
$(H\oplus N)^\perp$ is an even unimodular lattice of rank $18$ and signature
$(1,17)$, and hence
\[
(H\oplus N)^\perp\;\cong\;H\oplus E_8\oplus E_8\;=\;M.
\]
\end{lemma}

\begin{proof}
Write $S:=H\oplus N\subset\Lam$. Since $H$ and $N$ are nondegenerate and
mutually orthogonal, $S$ is nondegenerate of rank $4$ and signature
$(2,2)$; hence $S^\perp$ is nondegenerate of rank $22-4=18$ and signature
$(3,19)-(2,2)=(1,17)$. As a sublattice of the even lattice $\Lam$, the lattice
$S^\perp$ is even. By the anti-isometry of discriminant forms for nondegenerate
sublattices of a unimodular lattice \cite{Nik80}*{\S 1},
\[
q_{S^\perp}\;\cong\;-q_{S}\;=\;-\bigl(q_H\oplus q_N\bigr).
\]
But $H\cong U$ and $N\cong U$ are unimodular, so $q_S$ is trivial; hence
$q_{S^\perp}$ is trivial, i.e.\ $S^\perp$ is unimodular. Lemma
\ref{lem:evenunimodular} now gives $S^\perp\cong H\oplus E_8\oplus E_8$.
\end{proof}

\begin{corollary}\label{cor:mirrorperioddomain}
With $N$ as above, the saturated orthogonal complement of $M=H\oplus E_8\oplus
E_8$ in $\Lam$ is isometric to $H\oplus U$, a lattice of rank $4$ and
signature $(2,2)$. Consequently the period domain of $M$-quasipolarised K3
surfaces is a $2$-dimensional type~IV domain, and
$(L^\vee)^\vee\cong H$ after identifying $L^\vee\cong M$: the mirror
construction is involutive.
\end{corollary}

\begin{proof}
The saturated complement of $(H\oplus N)^\perp$ is $H\oplus N$ itself, of rank
$4$ and signature $(2,2)$. The rest follows from the definition of the period
domain in \S 3 and a second application of Lemma \ref{lem:mirrorlattice} with
the roles of the two sides exchanged (the complement of $H\oplus N$ in the
$M$-polarised picture is computed by the same argument).
\end{proof}

\begin{remark}
The isomorphism of Lemma \ref{lem:mirrorlattice} is independent of the
$O(\Lam)$-orbit of the cusp $N$: \emph{every} primitive isotropic plane
$N\cong U$ yields the same mirror lattice $M$. This is what makes the statement
of Theorem \ref{thm:main} uniform in the given degeneration. Had the cusp
lattice been of type $U(n)$ with $n\ge 2$, the discriminant form
$q_{S^\perp}\cong -q_{U(n)}$ would be nontrivial and the mirror lattice would
not be $M$; this is precisely the alternative excluded by Lemma
\ref{lem:divisibility} and Proposition \ref{prop:cusplattice}(ii).
\end{remark}

\section{The mirror surface and its elliptic fibration}\label{sec:fibration}

Let $Y$ be a very general $M$-quasipolarised K3 surface, with quasipolarisation
$j\colon M\hookrightarrow\NS(Y)$. By Lemma \ref{lem:verygeneralNS}:
\begin{equation}\label{eq:NSisM}
j\colon M\;\stackrel{\sim}{\longrightarrow}\;\NS(Y),\qquad \rho(Y)=18.
\end{equation}
Write $M=H\oplus E_8\oplus E_8$ and let $\e,\ff$ denote the standard isotropic
basis of the $H$-summand.

\begin{lemma}[Existence of the elliptic fibration]\label{lem:fibrationexists}
There exists an elliptic fibration $\pi\colon Y\to\Pn$ with a section, whose
fibre class is the primitive isotropic class $f:=j(\ff)$ (after possibly
replacing $j$ by $j\circ w$ for a Weyl isometry $w$ of $M$, which does not
affect any of the data in Theorem \ref{thm:main}). Moreover the class
$\sigma_0:=j(\e-\ff)$ satisfies $\sigma_0^2=-2$, $\sigma_0\cdot f=1$, is
effective, and is represented by a section of $\pi$.
\end{lemma}

\begin{proof}
By Saint-Donat's theorem \cite{SD74}*{Thm.~2.1} (see
\cite{Huy16}*{\S2.3}), a primitive isotropic class on a K3 surface is, up to
sign and up to the action of the Weyl group of $(−2)$-classes in $\NS(Y)$, nef;
and a nef primitive isotropic class $f$ with $f^2=0$ defines a base-point-free
elliptic pencil $\pi\colon Y\to\Pn$ unless it is the pullback of
$\OO_{\mathbb{P}^2}(1)$ under a double cover, an alternative excluded since such
a class has self-intersection $2$. Replacing $j$ by $j\circ w$ for a suitable
Weyl isometry $w$ (which changes neither the $M$-quasipolarisation in the sense
of \cite{Dol96} nor any of the lattice data below, since
$w(M)=M$), we may and do assume $f=j(\ff)$ is nef.

The class $\sigma_0=j(\e-\ff)$ has $\sigma_0^2=-2$ and $\sigma_0\cdot f=1$.
By Riemann--Roch on a K3 surface, $\chi(\sigma_0)=2+\sigma_0^2/2=1$, so
$h^0(\sigma_0)+h^0(-\sigma_0)\geq 1$ and one of $\pm\sigma_0$ is effective;
since $-\sigma_0\cdot f=-1<0$ while $f$ is nef, $-\sigma_0$ is not effective.
Hence $\sigma_0$ is effective, and any effective divisor $C\in|\sigma_0|$
satisfies $C\cdot f=1$, so exactly one irreducible component of $C$ is a curve
meeting the general fibre of $\pi$ transversally in one point: a section.
\end{proof}

\begin{lemma}[Mordell--Weil rank and the root budget]\label{lem:shiodatate}
Let $R_v\subset \NS(Y)$ denote the root lattice spanned by the classes of the
non-identity components of the fibre of $\pi$ over $v$, for each $v\in\Pn$.
Then:
\begin{enumerate}
\item[(a)] $R_v\subset E:=j(E_8\oplus E_8)=M\cap\langle f,\sigma_0\rangle^\perp$,
where $\langle f,\sigma_0\rangle\cong H$ (so that $E\cong E_8\oplus E_8$);
\item[(b)] the Shioda--Tate formula \cite{Shi72,Mir89} gives
$\operatorname{rank}\operatorname{MW}(\pi)+\sum_v\operatorname{rank}R_v=16$;
in particular $\sum_v\operatorname{rank}R_v\leq 16$.
\end{enumerate}
\end{lemma}

\begin{proof}
(a) The classes of fibre components are orthogonal to the fibre class $f$, and
the classes of non-identity components are moreover orthogonal to the section
class $\sigma_0$, since a section meets each fibre in a single point lying on
the identity component \cite{Mir89,Kod63}. Since
$\NS(Y)=M$ by \eqref{eq:NSisM} and $\langle f,\sigma_0\rangle=\langle
j(\ff),j(\e-\ff)\rangle$ is the $H$-summand of $M$, its orthogonal complement
inside $M$ is $j(E_8\oplus E_8)$. This is a copy of $E_8\oplus E_8$ by the same
discriminant argument as in Lemma \ref{lem:mirrorlattice} (even, negative
definite, unimodular of rank $16$).
(b) The Shioda--Tate formula states
$\rho(Y)=2+\operatorname{rank}\operatorname{MW}(\pi)+\sum_v\operatorname{rank}R_v$
\cite{Shi72} (see \cite{Mir89}*{Ch.~VII}); inserting $\rho(Y)=18$ gives the
claim.
\end{proof}

\begin{lemma}[Order of vanishing data]\label{lem:jorders}
Let $j_Y\colon\Pn\to\Pn$ be the functional $j$-invariant of the fibration
$\pi$, and let $e(F)$ denote the Euler characteristic of a fibre $F$.
Then:
\begin{enumerate}
\item[(a)] $\sum_v e(F_v)=e(Y)=24$ (Ogg's formula
\cite{Mir89,Kod63});
\item[(b)] if $\pi$ is non-isotrivial, $j_Y$ is a rational map of degree
$d=\sum_{I_n}n+\sum_{I_n^*}n\geq 1$, and the total order of vanishing of $j_Y$
at the points where $j_Y=0$ equals $d$;
\item[(c)] the orders of $j_Y$ at the fibres of each Kodaira type are
\[
\begin{array}{c|cccccccccc}
\text{type} & I_n & I_n^* & \mathrm{II} & \mathrm{III} & \mathrm{IV} &
I_0^* & \mathrm{IV}^* & \mathrm{III}^* & \mathrm{II}^*\\
\hline
\operatorname{ord} j_Y & -n & -n & 1 & 0 & 2 & 0 & 1 & 0 & 2
\end{array}
\]
as follows from Tate's algorithm via
$\operatorname{ord} j = 3\operatorname{ord} c_4-\operatorname{ord}\Delta$
\cite{Kod63,Mir89}.
\end{enumerate}
\end{lemma}

\begin{proposition}[Singular fibres of the mirror fibration]\label{prop:fibres}
Let $Y$ be a very general $M$-quasipolarised K3 surface and
$\pi\colon Y\to\Pn$ the fibration of Lemma \ref{lem:fibrationexists}. Then
$\pi$ has $\operatorname{rank}\operatorname{MW}=0$, and its singular fibres are
\emph{exactly} two of type $\mathrm{II}^*$ and four of type $\mathrm{I}_1$:
\[
\{\,F_v\;\}\;=\;\{\mathrm{II}^*,\,\mathrm{II}^*,\,I_1,\,I_1,\,I_1,\,I_1\,\}.
\]
\end{proposition}

\begin{proof}
\emph{Step 1 (a model family).} Consider Weierstrass equations
\begin{equation}\label{eq:weierstrass}
y^2=x^3+a(t)x+b(t),\qquad
a(t)=c_a(t-\alpha)^4(t-\beta)^4,\quad
b(t)=(t-\alpha)^5(t-\beta)^5\,b_2(t),
\end{equation}
where $t$ is a coordinate on $\Pn$, $b_2$ ranges over quadratics, and
$\alpha\neq\beta$. A direct computation via
$\Delta=-16(4a^3+27b^2)$ gives
\[
\Delta=(t-\alpha)^{10}(t-\beta)^{10}\,R(t),\quad
R(t)=4c_a^3(t-\alpha)^2(t-\beta)^2+27b_2(t)^2,
\]
so $\alpha,\beta$ carry fibres of type $\mathrm{II}^*$
($\operatorname{ord}c_4=4$, $\operatorname{ord}\Delta=10$
\cite{Kod63,Mir89}), and, for $c_a\neq0$ and $R$ with four simple roots
off $\{\alpha,\beta\}$, exactly four further singular fibres, all of type
$I_1$ (simple zero of $\Delta$, $c_4$ a unit there). The parameter count is:
$c_a$ ($1$), $\alpha,\beta$ ($2$), $b_2$ ($3$), i.e.\ $6$ parameters, against
the $4$-dimensional automorphism group ($1$ from the
$(x,y)\mapsto(\lambda^2x,\lambda^3y)$ scaling plus $3$ from
$\operatorname{PGL}_2$ acting on $t$), giving a
\emph{two}-dimensional family of fibrations with frame lattice containing
$H\oplus E_8\oplus E_8=M$. Since the moduli of $M$-quasipolarised K3 surfaces is
$2$-dimensional (Corollary \ref{cor:mirrorperioddomain}) and a $k$-dimensional
complete family of K3 surfaces has generic Picard rank at most $20-k$, the very
general member of this family has $\rho=18$ and $\NS=M$; it is very general.
Hence the configuration $(\mathrm{II}^*,\mathrm{II}^*,I_1^{\times4})$ occurs on
a nonempty, locally closed, two-dimensional --- hence dense --- stratum of the
irreducible two-dimensional moduli of $M$-quasipolarised K3 surfaces, on which
it is locally constant (Kodaira types are locally constant in families away
from collisions of singular values \cite{Mir89,Per85}).

\emph{Step 2 (alternatives with trivial Mordell--Weil rank).} Suppose
$\operatorname{rank}\operatorname{MW}=0$, so
$\sum_v\operatorname{rank}R_v=16$ by Lemma \ref{lem:shiodatate}(b), and
suppose $\pi$ is non-isotrivial. Since the fibre root lattices of Kodaira type
are $A_{n-1}$ ($I_n$), $D_{n+4}$ ($I_n^*$), $A_2,A_1$ ($\mathrm{IV}$,
$\mathrm{III}$), $D_4$ ($I_0^*$), $E_6,E_7,E_8$ ($\mathrm{IV}^*$,
$\mathrm{III}^*$, $\mathrm{II}^*$), and the total root rank is $16$, the
configuration is a decomposition of a full-rank sublattice of
$E_8\oplus E_8$ into Kodaira root systems. Combining
Lemma \ref{lem:jorders}(a),(b),(c): writing
$\Sigma e-\Sigma\operatorname{rank}=m+2a$ with $m$ the number of
multiplicative fibres and $a$ the number of additive fibres, we get
$24-16=8=m+2a$; and the zero-order budget $\sum\operatorname{ord}_0j_Y=d$
(from Lemma \ref{lem:jorders}(b),(c)) rules out, one by one, all decompositions
of $E_8\oplus E_8$ into Kodaira root systems (as classified via the maximal
root subsystems of $E_8$ \cite{Dyn57}) except two:
\[
\{\mathrm{II}^*,\mathrm{II}^*,I_1,I_1,I_1,I_1\}
\quad\text{and}\quad
\{\mathrm{II}^*,\mathrm{III}^*,\mathrm{III},I_1,I_1\},
\]
the latter carrying $2$-torsion in $\operatorname{MW}$. Indeed, e.g.\ the
decompositions $(E_8,E_7\!\oplus\!A_1)$, $(E_8,E_6\!\oplus\!A_2)$,
$(E_8,D_8)$, $(E_8,A_8)$, $(E_6\!\oplus\!A_2)^2$,
$(E_7\!\oplus\!A_1)^2$ all violate the zero/pole budget of
Lemma \ref{lem:jorders} (e.g.\ $(\mathrm{II}^*,\mathrm{IV}^*,\mathrm{IV},I_1^2)$
has $d=2$ but total zero order $2+1+2=5$; $(E_7\oplus A_1)^2$ gives $j\equiv
1728$, isotrivial; etc.).

\emph{Step 3 (the alternative strata are proper).} A Weierstrass normal form
with fibres $\mathrm{II}^*,\mathrm{III}^*,\mathrm{III}$ forces
$a(t)=(t-\alpha)^4(t-\beta)^3(a_0+a_1t)$ and
$b(t)=(t-\alpha)^5(t-\beta)^5b_2(t)$ together with a further condition cutting
out the $\mathrm{III}$ fibre; the parameter count is $2+2+3=7$ against the same
$4$-dimensional group plus one further condition, giving a
\emph{one}-dimensional family: a proper sublocus. The isotrivial alternative
$\{\mathrm{II}^*,\mathrm{II}^*,\mathrm{II},\mathrm{II}\}$ requires $a\equiv0$
(i.e.\ $j\equiv0$), a codimension-one condition, again proper. Finally, if
$\operatorname{rank}\operatorname{MW}>0$, the fibration's frame
($\langle f,\sigma_0\rangle\oplus\bigoplus_v R_v$ together with the
Mordell--Weil lattice \cite{Shi72}) still has rank $18$, but every such
configuration with frame $M$ other than
$(\mathrm{II}^*,\mathrm{II}^*,I_1^{\times4})$ occurs, by the classification of
fibre configurations on elliptic K3 surfaces due to Persson
\cite{Per85} (see \cite{Mir89}*{Ch.~X}), only on families of dimension strictly
less than $20-\operatorname{rk}(\text{frame})=2$. Since the moduli of
$M$-quasipolarised K3 surfaces is an irreducible $2$-dimensional period domain
quotient (Corollary \ref{cor:mirrorperioddomain}), all these alternative strata
are proper closed subsets, and a very general $Y$ avoids all of them, as well
as all Noether--Lefschetz divisors. Hence the configuration of
Step 1 prevails, $\operatorname{rank}\operatorname{MW}=0$, and the singular
fibres are as claimed.
\end{proof}

\begin{remark}
The conclusion $\operatorname{rank}\operatorname{MW}=0$ can also be seen
directly: the two $E_8$-summands of $M$ are then each the root lattice of a
single fibre (any proper root subsystem of $E_8$, such as
$E_7\oplus A_1$, sits in $E_8$ with index $>1$ \cite{Dyn57}, and index $>1$
would produce Mordell--Weil torsion, forcing one of the proper sub loci above).
\end{remark}

\begin{proposition}[Allowability of the loop]\label{prop:loop}
Let $\gamma\subset\Pn\setminus\Sigma(\pi)$ be a simple loop based at a general
point, enclosing exactly one of the two $\mathrm{II}^*$ fibres. Then
$\gamma$ is an allowable loop in the sense of
\cite{dfk3slpms}*{Definition 4.11}: its monodromy acting on
$H^1$ of the smooth fibre has order $6$, and the induced monodromy on
$H^2(Y;\Z)$ is precisely of the type matched, in
\cite{dfk3slpms}*{\S\S 4--5}, with the Type~II monodromy of the degeneration
$\cX\to\Delta$.
\end{proposition}

\begin{proof}
By Kodaira's classification \cite{Kod63}*{Tables}, the local monodromy of a
fibre of type $\mathrm{II}^*$ is conjugate to a matrix of order $6$
(concretely $\left(\begin{smallmatrix}0&-1\\1&-1\end{smallmatrix}\right)$ up to
inversion), acting on $H^1$ of the smooth fibre. This is the additive,
order-$6$ monodromy which \cite{dfk3slpms}*{Definition 4.11, Remark 4.12}
declares to be the allowable class for loops; geometrically, filling in such a
loop produces exactly the Type~II (Kulikov) degeneration behaviour matched
with the Tyurin degeneration by the mirror construction
\cite{dfk3slpms}*{\S\S 4--5}. Since by Proposition \ref{prop:fibres} the loop
encloses a single fibre of type $\mathrm{II}^*$ and no other singular fibre,
$\gamma$ is allowable.
\end{proof}

\begin{proposition}[Compatibility]\label{prop:compatibility}
The $M$-quasipolarisation $j\colon M\to\NS(Y)$ is compatible with the fibration
$\pi\colon Y\to\Pn$ and the loop $\gamma$ in the sense of
\cite{dfk3slpms}*{Definition 4.21}.
\end{proposition}

\begin{proof}
By construction, the fibre class $f=j(\ff)$ and the section class
$\sigma_0=j(\e-\ff)$ lie in $j(M)$ and generate its $H$-summand; the two
$E_8$-summands $j(E_8\oplus E_8)$ are exactly the root lattices of the fibres
inside and outside $\gamma$ (Proposition \ref{prop:fibres}); and the monodromy
around $\gamma$ preserves $j(M)$, acting trivially on $f$, $\sigma_0$, and on
the $E_8$-summand corresponding to the fibre outside $\gamma$, and by a Weyl
isometry on the $E_8$-summand of the enclosed fibre. These are precisely the
conditions of \cite{dfk3slpms}*{Definition 4.21}, as spelled out in
\cite{dfk3slpms}*{\S 4}.
\end{proof}

\section{Assembly: proof of the main theorem}\label{sec:proof}

The engine of the proof is the central construction of
\cite{dfk3slpms}, which we now state in our notation, incorporating the
correction to \cite{dfk3slpms}*{Definition 3.20} noted in
Definition \ref{def:hpolariseddegeneration}.

\begin{theorem}[{Giovenzana--Thompson, \cite{dfk3slpms}*{\S 5}}]\label{thm:GT}
Let $L\subset\Lam$ be a primitively embedded even lattice and let
$\cX\to\Delta$ be a very general Tyurin degeneration of $L$-polarised K3
surfaces in the sense of the corrected \cite{dfk3slpms}*{Definition 3.20}. Let
$N\subset L^\perp$ be the cusp lattice of the degeneration and
$L^\vee:=(L\oplus N)^\perp\subset\Lam$ the mirror lattice. Then there exists an
$L^\vee$-quasipolarised K3 surface $Y$, which may be taken very general,
together with an elliptic fibration $\pi\colon Y\to\Pn$ and an allowable loop
$\gamma\subset\Pn$ (in the sense of \cite{dfk3slpms}*{Definition 4.11}), such
that the $L^\vee$-quasipolarisation of $Y$ is compatible with $\pi$ and
$\gamma$ (in the sense of \cite{dfk3slpms}*{Definition 4.21}) and such that
$\cX\to\Delta$ and $\pi\colon Y\to\Pn$ form a mirror pair in the sense of
\cite{dfk3slpms}*{Definition 5.1}.
\end{theorem}

\begin{proof}[Proof of Theorem \ref{thm:main}]
Apply Theorem \ref{thm:GT} with $L=H$. The hypotheses are satisfied by
assumption: $\cX\to\Delta$ is a Tyurin degeneration of K3 surfaces which is
$H$-polarised in the sense of the corrected Definition
\ref{def:hpolariseddegeneration} (conditions (1) and (2), including
$(A|_{V_i})^2>0$, are part of the statement of Theorem \ref{thm:main}), and it
is very general.

By Proposition \ref{prop:cusplattice}, the degeneration determines the cusp
lattice $N\subset H^\perp\cong 2H\oplus 2E_8$ with $N\cong U$ (here Lemma
\ref{lem:divisibility} is the essential input). By Lemma
\ref{lem:mirrorlattice}, the mirror lattice is
\[
L^\vee=(H\oplus N)^\perp\;\cong\;H\oplus E_8\oplus E_8\;=\;M.
\]
Theorem \ref{thm:GT} therefore produces a very general
$M$-quasipolarised K3 surface $Y$, an elliptic fibration
$\pi\colon Y\to\Pn$, and an allowable loop $\gamma$, compatible with the
$M$-quasipolarisation, such that $(\cX\to\Delta,\;\pi\colon Y\to\Pn,\gamma)$ is
a mirror pair in the sense of \cite{dfk3slpms}*{Definition 5.1}.

It remains to identify the fibration and the loop on an arbitrary very general
$M$-quasipolarised $Y$, as asserted in the statement of Theorem
\ref{thm:main}. By Lemma \ref{lem:verygeneralNS}, every very general
$M$-quasipolarised K3 surface has $\NS(Y)=j(M)$; by Lemma
\ref{lem:fibrationexists} it carries the elliptic fibration with fibre class
$j(\ff)$ and section class $j(\e-\ff)$; by Proposition \ref{prop:fibres} this
fibration has exactly two $\mathrm{II}^*$ fibres and four $I_1$ fibres, in
agreement with the component lattice $R(X_0)\cong E_8\oplus E_8$ of
Proposition \ref{prop:componentlattice} (each $E_8$ being the root lattice of
one $\mathrm{II}^*$ fibre, one enclosed by $\gamma$ and one outside); by
Proposition \ref{prop:loop} a loop $\gamma$ enclosing exactly one
$\mathrm{II}^*$ fibre is allowable; and by Proposition \ref{prop:compatibility}
the quasipolarisation is compatible with $(\pi,\gamma)$. Finally, the data of
the pair $(Y,\pi,\gamma)$ that enter the mirror-pair condition of
\cite{dfk3slpms}*{Definition 5.1} --- the mirror lattice $M$, the
degeneration lattice $E_8\oplus E_8$, the cusp/loop lattice $N\cong U$, and the
Type~II monodromy data --- depend on the very general $M$-quasipolarised
surface $Y$ only through $\NS(Y)=j(M)$, by
Lemmas \ref{lem:fibrationexists}, \ref{lem:shiodatate} and Propositions
\ref{prop:fibres}--\ref{prop:compatibility}; hence they coincide with the data
produced by Theorem \ref{thm:GT}, and the very general $Y$ of the statement,
equipped with the fibration of \S\ref{sec:fibration} and a loop around one of
its $\mathrm{II}^*$ fibres, yields a mirror pair with $\cX\to\Delta$. This
completes the proof.
\end{proof}

\begin{remark}
If one reads \cite{dfk3slpms}*{Definition 5.1} as including, beyond the
lattice and monodromy data, a period-level matching between the limit mixed
Hodge structure of the degeneration and the Hodge-theoretic data of the
fibration around $\gamma$, then $Y$ should be taken to be the very general
surface produced by the construction of \cite{dfk3slpms}*{\S 5}, which is very
general precisely because the degeneration $\cX\to\Delta$ is assumed very
general (the Dolgachev correspondence \cite{Dol96,dfk3slpms} exchanges
Noether--Lefschetz loci on the two sides). The statement of Theorem
\ref{thm:main} is valid in either reading.
\end{remark}

\section{Concluding remarks}

\begin{remark}[Role of the corrected hypothesis]
The added condition $(A|_{V_i})^2>0$ in Definition
\ref{def:hpolariseddegeneration}(2) is essential twice: it ensures, via the
Hodge index theorem on each component, that the degeneration lattice
$R(X_0)$ is a sum of negative definite root lattices (Proposition
\ref{prop:componentlattice}), so that it can be matched against the root
lattices of the singular fibres of the mirror fibration; and it underlies the
nefness/isotropy analysis that forces the components of a very general
$H$-polarised Tyurin degeneration to be rational elliptic surfaces.
\end{remark}

\begin{remark}[Symmetry of the construction]
By Corollary \ref{cor:mirrorperioddomain}, the mirror of an
$M$-quasipolarised K3 surface, degenerating in Type~II at its unique cusp, is
an $H$-polarised K3 surface; the configuration
$\{\mathrm{II}^*,\mathrm{II}^*,I_1^{\times4}\$ of Proposition \ref{prop:fibres}
is the mirror avatar of the lattice $E_8\oplus E_8$ of
Proposition \ref{prop:componentlattice}, in exact analogy with the
$A_{17}$ example of \cite{dfk3slpms}*{\S\S 3--5}, where the lattice
$A_{17}$ arising from the components $\mathbb{P}^2\cup
\operatorname{Bl}_{18}\mathbb{P}^2$ corresponds to a mirror fibration with an
$I_{18}$ fibre. The two examples together illustrate that the fibration side of
the mirror pair is \emph{not} determined by the cusp lattice alone, but by the
full degeneration lattice $R(X_0)$.
\end{remark}

\begin{thebibliography}{ABE22}

\bibitem[ABE22]{cmek3sspt} V.~Alexeev, A.~Brunyate, and P.~Engel,
\emph{Compactifications of moduli of elliptic {K3} surfaces: Stable pair and
toroidal}, Geom. Topol. \textbf{26} (2022), no.~8, 3525--3588.

\bibitem[Dol96]{Dol96} I.~Dolgachev,
\emph{Mirror symmetry for lattice-polarized {K3} surfaces},
J. Math. Sci. \textbf{81} (1996), no.~3, 2599--2630.

\bibitem[DPT24]{DPT24} C.~F. Doran, J.~Prebble, and A.~Thompson,
\emph{Normal forms and {T}yurin degenerations of {K3} surfaces polarised by a
rank 18 lattice}, Preprint, arXiv:2311.10394.

\bibitem[Dyn57]{Dyn57} E.~B. Dynkin,
\emph{Semisimple subalgebras of semisimple {L}ie algebras},
Amer. Math. Soc. Transl. (2) \textbf{6} (1957), 111--243.

\bibitem[Fri83]{Fri83} R.~Friedman,
\emph{Global smoothings of varieties with normal crossings},
Ann. of Math. (2) \textbf{118} (1983), no.~1, 75--114.

\bibitem[GT24]{dfk3slpms} L.~Giovenzana and A.~Thompson,
\emph{Degenerations and fibrations of {K3} surfaces: Lattice polarisations and
mirror symmetry}, Preprint, May 2024, arXiv:2405.12009.

\bibitem[GHS07]{GHS07} V.~Gritsenko, K.~Hulek, and G.~K. Sankaran,
\emph{The {K}odaira dimension of the moduli of {K3} surfaces},
Invent. Math. \textbf{169} (2007), no.~3, 459--508.

\bibitem[Har77]{Har77} R.~Hartshorne, \emph{Algebraic geometry},
Graduate Texts in Mathematics, vol. 52, Springer-Verlag, New York, 1977.

\bibitem[Huy16]{Huy16} D.~Huybrechts, \emph{Lectures on {K3} surfaces},
Cambridge Studies in Advanced Mathematics, vol. 158, Cambridge University
Press, Cambridge, 2016.

\bibitem[Kod63]{Kod63} K.~Kodaira,
\emph{On compact analytic surfaces: {II}, {III}},
Ann. of Math. (2) \textbf{77} (1963), 563--626; \textbf{78} (1963), 1--40.

\bibitem[Kul77]{Kul77} V.~S. Kulikov,
\emph{Degenerations of {K3} surfaces and {E}nriques surfaces},
Math. USSR Izv. \textbf{11} (1977), 957--989.

\bibitem[Mir89]{Mir89} R.~Miranda, \emph{The basic theory of elliptic surfaces},
Dottorato di Ricerca in Matematica, ETS Editrice, Pisa, 1989.

\bibitem[Nik80]{Nik80} V.~V. Nikulin,
\emph{Integer symmetric bilinear forms and some of their geometric
applications}, Math. USSR Izv. \textbf{14} (1980), 103--167.

\bibitem[Per85]{Per85} U.~Persson,
\emph{Configurations of {K}odaira fibers on elliptic {R}iemann surfaces},
in: Algebraic geometry (Sitges, 1983), Lecture Notes in Math., vol. 1124,
Springer, Berlin, 1985, pp. 334--369.

\bibitem[PP81]{PP81} U.~Persson and H.~Pinkham,
\emph{Degenerations of surfaces with anticanonical bundles},
J. Differential Geom. \textbf{13} (1978/81), no.~2, 209--238.

\bibitem[PS71]{PS71} I.~I. Pjatecki{\u\i}-\v{S}apiro and I.~R. \v{S}afarevi\v{c},
\emph{A {T}orelli theorem for algebraic surfaces of type {K3}},
Math. USSR Izv. \textbf{5} (1971), 547--588.

\bibitem[SB83]{SB83} N.~I. Shepherd-Barron,
\emph{Degenerations of {K3} surfaces}, 1983, manuscript cited in
\cite{DPT24}.

\bibitem[SD74]{SD74} B.~Saint-Donat,
\emph{Projective models of {K3} surfaces},
Amer. J. Math. \textbf{96} (1974), 602--639.

\bibitem[Ser70]{Ser70} J.-P. Serre, \emph{A course in arithmetic},
Graduate Texts in Mathematics, vol. 7, Springer-Verlag, New York, 1973.

\bibitem[Shi72]{Shi72} T.~Shioda,
\emph{On elliptic surfaces in characteristic zero},
in: Commentationes Mathematicae Universitatis Sancti Pauli
(vol.~21, 1972/73); see also \emph{An explicit algorithm for computing the
Picard number of certain elliptic surfaces}, Amer. J. Math. \textbf{108} (1986).

\bibitem[Ste77]{Ste77} J.~H.~M. Steenbrink,
\emph{Limits of {H}odge structures},
Invent. Math. \textbf{31} (1976/77), no.~3, 229--257.

\bibitem[SZ01]{SZ01} I.~Shimada and D.-Q. Zhang,
\emph{Classification of extremal elliptic {K3} surfaces and the degree of the
curve of intersection of two quadrics}, Michigan Math. J. \textbf{50} (2002).

\end{thebibliography}

\end{document}
```